\documentclass[aspectratio=169]{beamer}

\usepackage[utf8]{inputenc}
\usepackage[T1]{fontenc}
\usepackage[english]{babel}
\usepackage{graphicx}
\usepackage{qrcode}
\usepackage{hyperref}
\usepackage{fontawesome5}
\usepackage{xcolor}
\usepackage{ragged2e}

\usetheme{Madrid}
\usecolortheme{default}

% Убираем навигационные символы
\setbeamertemplate{navigation symbols}{}

\title{\textbf{Calorie Counter}}
\subtitle{A Telegram bot for personalised calorie, protein, fat, and carb tracking}
\author{Iuliana Giliazutdinova \\ \textit{\small i.giliazutdinova@innopolis.university}}
\institute{Innopolis University \\ Group B25-DSAI-02}
\date{\today}

\begin{document}

% ==================== СЛАЙД 1: Титульный ====================
\begin{frame}[plain]
    \titlepage
\end{frame}

% ==================== СЛАЙД 2: Контекст ====================
\begin{frame}{Context}
    \vspace{0.3cm}
    
    \begin{block}{\faUser{} End-User}
        People who track their daily nutrition intake.
    \end{block}

    \vspace{0.3cm}
    
    \begin{block}{\faExclamationTriangle{} What Problem Does It Solve?}
        People want to track daily nutrition, but existing apps are cluttered, push premium paywalls, use generic norms, and require too many taps to log a meal. For users with eating disorders, calorie-focused interfaces can be triggering.
    \end{block}

    \vspace{0.3cm}
    
    \begin{alertblock}{\faLightbulb{} The Product Idea in One Sentence}
        Our bot solves this: one message \texttt{200/30/15/45} logs calories, protein, fat, and carbs instantly, and shows progress toward a personalised daily norm calculated from the user's body, activity, and goals.
    \end{alertblock}
\end{frame}

% ==================== СЛАЙД 3: Реализация ====================
\begin{frame}{Implementation}
    \begin{columns}[T]
        \begin{column}{0.5\textwidth}
            \begin{block}{\faCode{} How It Was Built}
                \begin{itemize}
                    \small
                    \item \textbf{Backend:} Python + aiogram 3.x
                    \item \textbf{Database:} PostgreSQL (Railway)
                    \item \textbf{Client:} Telegram Bot API
                    \item \textbf{Deploy:} Railway.app
                \end{itemize}
            \end{block}

            \begin{block}{\faUsers{} Stakeholders}
                \begin{itemize}
                    \small
                    \item End-users seeking a quick, stress-free way to log daily nutrition
                    \item Fitness enthusiasts wanting personalised norms tailored to their body and goals
                \end{itemize}
            \end{block}
        \end{column}
        \begin{column}{0.5\textwidth}
            \begin{block}{\faCogs{} Version 1: Core Tracker}
                \begin{itemize}
                    \small
                    \item Personalised daily norm calculation (gender, weight, height, age, activity level, goal)
                    \item One-message food logging (calories/protein/fat/carbs)
                    \item Daily summary with progress tracking
                    \item Backend (PostgreSQL) + Telegram client
                \end{itemize}
                \smallskip
                \textbf{Outcome:} A working MVP that delivers the core value: log intake, instantly view progress toward personalised goals.
            \end{block}
        \end{column}
    \end{columns}
\end{frame}

% ==================== СЛАЙД 4: Демо ====================
\begin{frame}{Demo: Version 2}
    \vspace{0.3cm}
    
    \begin{columns}[T]
        \begin{column}{0.55\textwidth}
            \begin{block}{\faRocket{} Version 2: Analytics \& Refinement}
                \begin{itemize}
                    \small
                    \item Interactive calendar with daily navigation and history view
                    \item Weekly/monthly statistics with consolidated totals
                    \item Advanced logging options (per-100g auto-calculation, entry correction/deletion)
                    \item Deployment hardening, edge-case handling, and integration of TA feedback
                \end{itemize}
                \smallskip
                \textbf{Outcome:} A production-ready app with a complete tracking lifecycle, historical analytics, and a polished UX for sustained use.
            \end{block}
        \end{column}
        \begin{column}{0.45\textwidth}
            \begin{center}
                \Huge \faVideoCamera{}
                \vspace{0.2cm}
                
                \textbf{Video Demonstration}
                \vspace{0.3cm}
                
                \href{https://www.youtube.com/watch?v=YOUR_VIDEO_ID}{
                    \color{blue}\underline{Watch the demo (under 2 min)}
                }
                
                \vspace{0.3cm}
                \small
                \textit{The video shows the full user journey from registration to viewing weekly stats.}
            \end{center}
        \end{column}
    \end{columns}
\end{frame}

% ==================== СЛАЙД 5: Ссылки ====================
\begin{frame}{Links}
    \vspace{0.3cm}
    
    \begin{columns}[T]
        \begin{column}{0.5\textwidth}
            \begin{center}
                \faGithub{} \textbf{GitHub Repository}
                \vspace{0.3cm}
                
                \href{https://github.com/kaftanovaa/kbzhyshka}{
                    \color{blue}\underline{github.com/kaftanovaa/kbzhyshka}
                }
                
                \vspace{0.3cm}
                
                \qrcode[height=3.5cm]{https://github.com/kaftanovaa/kbzhyshka}
            \end{center}
        \end{column}

        \begin{column}{0.5\textwidth}
            \begin{center}
                \faRocket{} \textbf{Deployed Product}
                \vspace{0.3cm}
                
                \href{https://t.me/kbzhyshka_bot}{
                    \color{blue}\underline{@kbzhyshka\_bot}
                }
                
                \vspace{0.3cm}
                
                \qrcode[height=3.5cm]{https://t.me/kbzhyshka_bot}
            \end{center}
        \end{column}
    \end{columns}
\end{frame}

\end{document}
